\clearpage
\section{Receiving the vineyard and bearing its fruit}
\zlabel{triptych:concise:commentary:start}
\subsection{What does the giver's priority change?}
The vineyard's first fact is that it has been planted. Chrysostom notices how much Matthew's householder does before the tenants arrive: the land is enclosed, a winepress dug, a tower built. He relates these provisions to God's gifts to Israel, including deliverance, law, temple and altar. The tenants inherit a cared-for world. Their task is real, but it begins within a gift. The requested fruit is not an arbitrary addition to their tenancy; it is the purpose for which the entrusted work exists.\footnote{Chrysostom, \work{Homilies on Matthew} 68.1, NPNF1-10.}

Augustine approaches the same priority by asking what Christ found in the disciples he chose. If prior goodness explained his choice, their turning toward him would have come first. John 15:16 places Christ's choosing first. Augustine therefore reads mercy as making the goodness it does not simply discover. Yet the disciples are chosen to go and bear fruit. Their dependence does not end at election, any more than a branch ceases to receive life once it begins bearing. Christ remains the way in which they go.\footnote{Augustine, \work{Tractates on John} 86.2--3, NPNF1-7. The appointed acclamation adapts John 15:16; this argument also uses its full biblical context.}

The two arguments expose different mistakes. Chrysostom challenges possession without accountability: the tenants treat the giver's provisions as their own unrestricted property. Augustine challenges goodness imagined as the cause that made Christ's first gift due. Both mistakes make the receiver independent of the giver. The practical answer is neither to neglect works nor to boast of them. Works become the response a gift enables and requires. The Collect's abundance beyond merit gives that response a fitting beginning in petition.

The wounded vineyard brings another need to this account. It may have lost the confidence that it can answer at all. The psalm does not tell it to produce an adequate recovery before asking God to return. It asks for quickening so that invocation becomes possible. The first reading of the formulary hears continued dependence here; the second hears hope for a damaged people. A person who needs help to act and a community that needs help to live can make the same petition for different, compatible reasons.

Chrysostom's Philippians homily adds an accent Augustine's account of election does not supply. At its close he urges the hearers to make a beginning and stop fighting. This is a summons to moral action within his earlier account of divine beneficence toward enemies. Augustine is answering what causes Christ's choice; Chrysostom is urging people already addressed by the apostolic exhortation to practice peace. Their common demand is active gratitude, while their distinct questions remain. Neither passage alone supplies a complete account of every relation between grace and initiative.\footnote{Chrysostom, \work{Homilies on Philippians} 14, exposition of 4:7 and concluding exhortation, NPNF1-13.}

The Entrance helps prevent dependence from becoming merely an abstract order of causes. Mordecai invokes the Creator amid a threat to his people. The surrounding Esther narrative joins prayer to action under danger. Sovereignty belongs to God, not to the frightened petitioner or the grasping tenant. Trust can therefore sustain action without promising the petitioner immediate control. Nothing in that danger establishes that the threatened people deserved destruction.

\subsection{Who is accused, and what would count as fruit?}
Isaiah's song leads its audience into a judgment. A carefully prepared vineyard should yield good grapes; the hearers can recognize its failure before the poem identifies the planting with Israel and Judah. The conclusion makes the fruit concrete: justice has given way to an outcry. The woes later in the chapter describe acquisitiveness and corrupted judgment. Their context sharpens the appointed passage's accusation. The deficiency is a wrong relation to other people, not simply an insufficient quantity of achievement.

Chrysostom compares this song with Matthew and distinguishes their subjects. Isaiah accuses the vineyard; Jesus accuses the husbandmen responsible for it. Matthew's tenants answer successive messengers with violence and calculate that killing the heir will secure his inheritance. Their wrongdoing is not merely disappointing work. They attempt to abolish the relationship that makes their stewardship possible. Repeated sending shows the owner's longsuffering and gives repeated occasions to respond; it does not make their answer less consequential.\footnote{Chrysostom, \work{Homilies on Matthew} 68.1--2. Matthew 21:45 identifies the chief priests and Pharisees as recognizing themselves in the parable; that verse is context beyond the appointed ending at 21:43.}

The owner's expectation that the son will be respected creates a difficulty: does the parable depict God as unaware of the outcome? Chrysostom reads it as naming the response that ought to have occurred. The tenants' obligation is genuine even though divine knowledge is not deficient. His account preserves both patience and culpability. The sending of the son reveals generosity at the point where the tenants reveal most completely their desire to seize.

Augustine gives the fruit a positive name by reading John 15:16 with the command that follows. Mutual love explains what lasting fruit means. Patience, kindness and peace take their order from charity, rather than standing as disconnected accomplishments. Beside Isaiah, this identification prevents charity from becoming benevolent sentiment detached from justice. The neighbor's good matters; the wronged person's outcry cannot be offset by a record of impressive religious activity.\footnote{Augustine, \work{Tractates on John} 87.1.}

The restoration reading asks what this criterion means for common life. A damaged community may want its former strength back while resisting the exposure of old wrongs. Isaiah makes that an inadequate hope. More secure walls would not repair the injustice within them. Augustine's moral exposition of the psalm accordingly turns to desires and fears: attraction to gain and dread of loss can each move a person to sin. Holy love and fear redirect those powers. Renewal reaches the motives by which members defend what harms their neighbors.\footnote{Augustine, \work{Exposition of Psalm} 79, §10, NPNF1-8, English heading Psalm 80.}

Judgment also remains directed toward new stewards. The people receiving the kingdom in Matthew 21:43 are characterized by fruitfulness, not by an entitlement to repeat the tenants' possessiveness. Chrysostom's harsh anti-Jewish rhetoric exceeds his more precise identification of the rulers; the same exposition speaks of believing Jews and Gentiles made one. The parable does not assign hereditary guilt to all Jewish people. Its criterion also challenges the Church's own conduct. Belonging to the vineyard cannot serve as proof that its owner receives his due.

\subsection{What survives when the vineyard is wounded?}
The psalm remembers the planting before it describes the broken hedge. The vine was brought from Egypt, given room and allowed to spread. Its present devastation is not imaginary, but neither is the gift preceding it. Memory addresses the planter concerning his own work. Asking him to look, visit and give life assumes that the wound need not have the last word. The accused vineyard of Isaiah and the praying vine of the psalm can belong to one people's story without making every sufferer equally culpable.

Augustine follows that identity through Christ. At the plea that God perfect what his right hand planted, he rejects the suggestion of an entirely other vineyard. The seed of Abraham and the blessing of the nations find their fulfillment in Christ; Gentiles enter as grafted participants in a root preceding them. Christ is the foundation of the perfected planting. Renewal is therefore not the disappearance of Israel's history into an unrelated origin story.\footnote{Augustine, Psalm 79, §§6,9.}

Bellarmine supplies a distinct, concrete argument for the same continuity. He contrasts reform with total destruction and names the Jewish apostles and converts who respond to Peter, alongside Gentile grafting. These people make the renewed vineyard's continuity visible. His claim is not merely that one agricultural image can be applied to different communities. Christ perfects and governs the planting, and the nations enter a promise they did not originate.\footnote{Bellarmine, \work{A Commentary on the Book of Psalms}, Psalm 79, at vv. 14--18, trans. John O'Sullivan (1866), abridged English witness.}

The Gospel's rejected Son belongs at the center of this answer. Augustine himself recalls the heir cast out of the vineyard in his psalm exposition: rejection does not prevent the heir from possessing it. Matthew's rejected stone likewise becomes the cornerstone. The stewards intend to settle the future by eliminating the Son, but their judgment cannot determine his place. Restoration is founded upon him, not upon a community's capacity to preserve its own unbroken reputation.\footnote{Augustine, Psalm 79, §7; Matthew 21:38--43.}

Agreement on Christ does not eliminate the commentators' historical differences. Augustine identifies the river as Jordan and relates the ravager to Roman power. Bellarmine identifies Euphrates and discusses Assyrian and Babylonian devastation. Those are different geographical and historical readings. Neither establishes the psalm's composition date, and their theological convergence does not turn the rival identifications into one reconstruction. Bellarmine's own assertion of agreement about Christ remains his assertion; it is not evidence that every interpreter agrees.\footnote{Augustine, Psalm 79, §§6--8; Bellarmine, Psalm 79, at vv. 10--13. Both adopted English psalm witnesses are abridged.}

This difference also clarifies the senses of Scripture. The literal vine is the people brought from Egypt, lamenting devastation. Its perfection in Christ is the allegorical fulfillment Augustine and Bellarmine expound. The gift reading emphasizes dependence upon Christ as the source of fruitful charity; the restoration reading emphasizes the continuity of the people gathered in him. Their moral consequences meet in gratitude and repentance, but renewal adds a particular demand: members must repair a common life, not merely improve separate private records.

\subsection{How does peace become the work of a people?}
Paul's exhortation addresses relationships as well as thoughts. Philippians 4 names Euodia and Syntyche before the appointed verses and speaks of the community's material care afterward. Thanksgiving and worthy attention stand within that common life. The final command reaches what the hearers have learned, received, heard and seen. Apostolic teaching is embodied in an example to practice. Peace is not an escape from the neighbor into a protected interior space.

Chrysostom begins by asking how Christian joy can coexist with mourning. He allows grief over sin and rejoicing in Christ together. Later he explains the peace beyond understanding by God's reconciliation of enemies through the Son. Its source is divine generosity, not a human technique for producing pleasant feelings. His opening acknowledgment of affliction controls his later strong exhortation about virtue: neither Paul nor his commentator promises immunity from painful events.\footnote{Chrysostom, \work{Homilies on Philippians} 14, opening and exposition of 4:6--9.}

For the gift reading, thanksgiving recalls what petitioners already receive. It resists the idea that God becomes valuable only when the next request succeeds. Attention to truth and justice then prepares action, while the apostolic example shows what action looks like. Chrysostom distinguishes doing what deserves praise from doing it in order to win praise. The inward direction of an apparently good deed matters. A vineyard can seem successful while the desire governing it remains the tenants' wish to possess.

For the restoration reading, the same teaching identifies habits that can heal or perpetuate division. Chrysostom names revenge, remembered injuries, envy, covetousness and vainglory. Refusing retaliation is not merely maintaining a quiet surface. God's generosity toward enemies challenges the rule that personal grievance should determine every response. Augustine similarly distinguishes love for the creature God made from rejection of the vice harming it. Love of an enemy seeks a person's good while opposing the evil.\footnote{Chrysostom, \work{Homilies on Philippians} 14, concluding exhortation; Augustine, \work{Tractates on John} 87.4.}

Truth and justice qualify this reconciliation from within Paul's own list. Neither a silence imposed upon the injured nor agreement around a false account would fulfill it. Isaiah's outcry remains audible beside the promise of peace. The community seeks an honest common life in which correction can be received and a neighbor's good pursued. The gift reading asks whether gratitude has become action; the restoration reading asks whether such action reaches the relations that have been damaged.

The Epistle has its own place in the semi-continuous apostolic course. Its movement from petition to practiced peace is richer than an instruction to attach the vineyard image to every sentence. Its contribution is precisely that it gives another concrete account of Christian action: thankful prayer, disciplined attention, an example received, and deeds performed. The vineyard images bring urgency to that action without exhausting what Paul says.

\subsection{Does Communion complete the response or make it possible?}
Aquinas begins his account of Eucharistic grace with Christ, whom the sacrament contains, and the Passion it represents. Only then does he explain nourishment: spiritual life is sustained, increased, restored and delighted. The causality comes from Christ; bread is not merely an illustration of people helping one another. The many grains and grapes also signify unity. Life received from one source binds the recipients in a common charity.\footnote{Aquinas, \work{Summa theologiae} III, q. 79, a. 1, body, historical English by the Fathers of the English Dominican Province.}

The two whole-formulary readings draw different consequences from that order. The gift reading finds the enabling source of the response. Aquinas says charity is aroused to action: reception is not a substitute for bearing fruit but nourishment for it. The restoration reading finds more than an improved arrangement of independent individuals. Christ gives the unity into which the many are incorporated. Better conduct expresses that gift; it does not manufacture a different Christ for each member.

Chrysostom develops the latter point directly from First Corinthians 10:16--17. Many grains become one bread; communicants receive the same Christ and belong both to him and to one another. He immediately tests the claim against loveless conduct. Why should people fed by the same gift refuse the same love? His sacramental affirmation makes division more serious. Communion is not a certificate that recipients already behave well; it reveals the common life their failure contradicts.\footnote{Chrysostom, \work{Homilies on First Corinthians} 24.4, NPNF1-12.}

That language is explicit in the Corinthians Communion alternative. The Missal adaptation includes chalice language whose biblical context is verse 16, so a bare quotation of verse 17 would not reproduce the entire antiphon. The other appointed alternative, Lamentations 3:25, gives a different accent to the same sacramental moment. Seeking and waiting occur within a chapter that has named lost peace and failed strength. Hope does not require pain to have disappeared before the communicant can approach.

Lamentations also rejects oppression and corrupted justice. Its patience cannot mean consenting to everything the powerful inflict. In the gift reading, waiting protects fruitfulness from being measured solely by visible results. In the restoration reading, it permits a wounded people to ask for life before the shape of repair is clear. The two antiphons are alternatives, not consecutive steps. The explicit one-body image is therefore not the compulsory verbal conclusion of this Sunday's Communion.\footnote{Lamentations 3:17--36,40--58; appointed alternative 3:25.}

The prayers around Communion give this reception a direction. The offerings prayer joins obedient worship with redemption's sanctifying action; its exact approved English remains uncollated, and the analysis rests on the identified Latin witness. The final prayer asks that nourishment shape the life of the receiver. Leo's Eucharistic teaching supplies a corresponding picture: those receiving Christ bear the crucified and risen Lord in body and spirit. His argument concerns participation in the Eucharist; the baptismal discussion belongs to the preceding section.\footnote{Leo the Great, Sermon 63, \work{On the Passion XII}, VII and note 1063, trans. Charles Lett Feltoe, NPNF2-12; the preceding VI discusses baptism.}

\subsection{What fulfillment can neither success nor ruin measure?}
Augustine closes his psalm exposition by asking what created gift could deserve greater love than its Creator. Wealth, health and honor are also possessed by people who do not worship God. If such things became the final reason for worship, another person's prosperity could appear to disprove God's goodness. Augustine instead makes God himself the reward. Created beauty invites love of its Maker; it cannot take his place.\footnote{Augustine, Psalm 79, §11.}

The gift reading consequently carries charity beyond a successful season. Lasting fruit is not exhausted by completed work, public esteem or recovered comfort. The restoration reading receives a related correction from Bellarmine's distinction between grace now and glory hereafter. A community may be renewed in charity while still bearing losses and scars. Neither its present strength nor its present wound measures the fullness of its promise.

Aquinas gives this hope a bodily reach. The body's members serve justice now and share future incorruption. The moral and anagogical senses therefore remain connected: embodied charity belongs to the present response, resurrection to its promised consummation. Yet they are not interchangeable. Present improvement is not glory already possessed, and future hope does not excuse neglect of present justice.\footnote{Aquinas, III, q. 79, a. 1, reply 3.}

The two readings finally answer different anxieties. The person trying to purchase divine favor hears that Christ's gift makes fruitful action possible. The people fearing that its wound is its whole identity hears that the planter can perfect what he planted. Both receive a summons as well as consolation. The vineyard must bear fruit; the renewed people must live its communion. Both remain dependent upon the rejected Son who is their cornerstone.
